\section{Proofs}
\label{app:proofs}

This appendix contains the complete finite-dimensional arguments. No statement below
depends on an optimiser, a conjecture, or a numerical plot.

\subsection{Magnitudes and the orthogonal splitting at the root}

For a subtree rooted at $v$, let $k_v$ be its number of internal vertices and $L_v$ the
product of the reduced leaf norms below it.

\begin{lemma}[Subtree magnitudes]
\label{lem:magnitudes}
If $\norm{\mu_u}_\op\le M$ at every internal vertex $u$ of the subtree, then
\[
 \norm{F_v}\le M^{k_v}L_v,
 \qquad
 \norm{R_v}\le M^{k_v}L_v.
\]
\end{lemma}

\begin{proof}
At a leaf, $F_v=R_v=Q_{\tau(v)}z_v$ and $Q_{\tau(v)}$ is an isometry, so both norms equal
$\norm{z_v}=L_v$, and $k_v=0$. Suppose the claim holds for the children
$c_1,\ldots,c_a$ of an internal vertex $v$. Multilinearity and the definition of the
operator norm give
\[
 \norm{F_v}
 \le M\prod_i\norm{F_{c_i}}
 \le M^{1+\sum_ik_{c_i}}\prod_iL_{c_i}
 =M^{k_v}L_v.
\]
The same argument bounds $\norm{R_v}=\norm{P_v\mu_v(R_{c_1},\ldots,R_{c_a})}$, because an
orthogonal projector is a contraction. Induction proves both claims.
\end{proof}

\begin{proof}[Proof of Proposition~\ref{prop:root}]
At the root write $P=QQ^*$, $F=F_\varrho$ and $R=R_\varrho$. If $\varrho$ is internal then
$R=P\mu_\varrho(\cdots)\in\ran P$; if $T$ consists of a single leaf then $R=Qz\in\ran P$
as well. In either case
\[
 F-R=(PF-R)+(I-P)F,
\]
where the first summand lies in $\ran P$ and the second in $(\ran P)^\perp$. The
Pythagorean theorem gives the first identity. For the second, $Q^*(I-P)=Q^*-Q^*QQ^*=0$,
so
\[
 Q^*F-Q^*R=Q^*(PF-R),
\]
and the restriction of $Q^*$ to $\ran P$ is an isometry onto $W_{\tau(\varrho)}$, because
$Q^*Q=I$. Hence $\norm{Q^*F-Q^*R}=\norm{PF-R}$.
\end{proof}

Both steps use $P^*=P$. Example~\ref{ex:oblique} shows that the conclusion fails for an
idempotent that is not self-adjoint.

\subsection{The exact local decomposition}

\begin{proof}[Proof of Theorem~\ref{thm:subset}]
Since $F_i=R_i+D_i$, repeated use of multilinearity in each argument gives
\[
 \mu_v(F_1,\ldots,F_a)
 =\sum_{S\subseteq[a]}\mu_v(y_1^S,\ldots,y_a^S),
\]
a sum of $2^a$ terms. The term indexed by $S=\varnothing$ is
$\mu_v(R_1,\ldots,R_a)$, while $R_v=P_v\mu_v(R_1,\ldots,R_a)$. Subtracting,
\[
 D_v
 =F_v-R_v
 =(I-P_v)\mu_v(R_1,\ldots,R_a)
 +\sum_{\varnothing\ne S\subseteq[a]}\mu_v(y_1^S,\ldots,y_a^S).
\]
It remains to identify the first summand with $r_v(R_1,\ldots,R_a)$. By construction
$R_i\in\ran P_{\tau(c_i)}$ for every $i$, so $P_{\tau(c_i)}R_i=R_i$, and therefore
\[
 r_v(R_1,\ldots,R_a)
 =(I-P_v)\mu_v\bigl(P_{\tau(c_1)}R_1,\ldots,P_{\tau(c_a)}R_a\bigr)
 =(I-P_v)\mu_v(R_1,\ldots,R_a).
\]
No inequality has been used.
\end{proof}

\begin{remark}
\label{rem:innerprojectors}
The inner projectors in Definition~\ref{def:closure} are not decorative. The hypothesis
$\norm{r_v}_\op\le\rho$ of Theorem~\ref{thm:homogeneous} bounds the operator
$(I-P_v)\mu_v(P\,\cdot,\ldots,P\,\cdot)$, whose operator norm is in general strictly
smaller than that of $(I-P_v)\mu_v(\cdot,\ldots,\cdot)$. The step
$P_{\tau(c_i)}R_i=R_i$ above is exactly what permits the smaller quantity to be used.
\end{remark}

Applying $P_v$ to Theorem~\ref{thm:subset} removes the locally created residual exactly;
applying $I-P_v$ retains it. Splitting each $D_i$ orthogonally into
$D_i^{\parallel}+D_i^{\perp}$ produces the three states $\{R,\parallel,\perp\}$ used in
Section~\ref{sec:state}.

\subsection{Uniform coefficients}

\begin{proof}[Proof of Theorem~\ref{thm:homogeneous}]
We prove, by induction on the tree, that every subtree rooted at $v$ satisfies
\begin{equation}
\label{eq:induction}
 \norm{D_v}\le k_v\rho M^{k_v-1}L_v
 \quad (k_v\ge1),
 \qquad\text{and}\qquad
 \norm{D_v}=0
 \quad (k_v=0).
\end{equation}
The two cases are stated separately because $M^{k-1}$ is not defined at $k=0$ when $M=0$;
compare Remark~\ref{rem:leaf}.

\emph{Base case.} At a leaf, $F_v=R_v$, so $\norm{D_v}=0$ and $k_v=0$.

\emph{Inductive step.} Let $v$ be internal of arity $a$ with children $c_1,\ldots,c_a$,
so that $k_v=1+\sum_ik_{c_i}$. Telescoping the arguments of $\mu_v$ one at a time, in any
fixed order, and using multilinearity together with $\norm{\mu_v}_\op\le M$,
\begin{equation}
\label{eq:telescope}
 \norm{\mu_v(F_1,\ldots,F_a)-\mu_v(R_1,\ldots,R_a)}
 \le M\sum_{j=1}^{a}\norm{F_j-R_j}\prod_{i<j}\norm{R_i}\prod_{i>j}\norm{F_i}.
\end{equation}
Fix $j$. If $k_{c_j}=0$ then $\norm{F_j-R_j}=0$ and the $j$th summand vanishes, which
agrees with $k_{c_j}\rho M^{k_v-1}L_v=0$. If $k_{c_j}\ge1$, the inductive hypothesis and
Lemma~\ref{lem:magnitudes} bound the $j$th summand by
\[
 M\cdot k_{c_j}\rho M^{k_{c_j}-1}L_{c_j}\cdot\prod_{i\ne j}M^{k_{c_i}}L_{c_i}
 =k_{c_j}\rho M^{\sum_ik_{c_i}}L_v
 =k_{c_j}\rho M^{k_v-1}L_v.
\]
Summing over $j$ and using $\sum_jk_{c_j}=k_v-1$, the right-hand side
of~\eqref{eq:telescope} is at most $(k_v-1)\rho M^{k_v-1}L_v$.

By Theorem~\ref{thm:subset},
$D_v=r_v(R_1,\ldots,R_a)+\bigl[\mu_v(F_1,\ldots,F_a)-\mu_v(R_1,\ldots,R_a)\bigr]$, and
the locally created residual satisfies
\[
 \norm{r_v(R_1,\ldots,R_a)}
 \le\rho\prod_i\norm{R_i}
 \le\rho M^{k_v-1}L_v
\]
by Lemma~\ref{lem:magnitudes} and $\sum_ik_{c_i}=k_v-1$. Adding the two contributions
gives $\norm{D_v}\le k_v\rho M^{k_v-1}L_v$, which is~\eqref{eq:induction}.

Applying~\eqref{eq:induction} at the root gives
$E_T^{\amb}=\norm{D_\varrho}\le k\rho M^{k-1}L_T$.

For the projected error, $R=R_\varrho\in\ran P$ gives $PR=R$, hence
$PF-R=P(F-R)=PD_\varrho$. Since $Pr_\varrho=0$ by Definition~\ref{def:closure},
\[
 PD_\varrho
 =P\bigl[\mu_\varrho(F_1,\ldots,F_a)-\mu_\varrho(R_1,\ldots,R_a)\bigr],
\]
and $\norm{P}=1$ together with the bound just established for the right-hand side
of~\eqref{eq:telescope} gives $E_T^{\proj}\le(k-1)\rho M^{k-1}L_T$.
Proposition~\ref{prop:root} identifies $E_T^{\red}$ with $E_T^{\proj}$.

For the orthogonal error, $(I-P)R=0$ gives $(I-P)F=(I-P)D_\varrho$, and $\norm{I-P}=1$
gives $E_T^{\perp}\le\norm{D_\varrho}\le k\rho M^{k-1}L_T$.
\end{proof}

\begin{remark}
\label{rem:whatisused}
The proof uses multilinearity, the operator-norm hypotheses, contractivity of $P$ and of
$I-P$, and --- for the projected estimate only --- the identity $Pr_\varrho=0$, which
requires $P$ to be the orthogonal projector appearing in Definition~\ref{def:closure}.
Finite dimensionality is not used: the argument is valid for bounded multilinear maps
between Hilbert spaces of any dimension. It is stated finite-dimensionally because
everything else in this article is.
\end{remark}

\subsection{The ordering theorem}

Let $e_i,r_i,f_i,G_i\ge0$, put $w_i=G_ie_i$ and $d_i=f_i-r_i$, and for a permutation
$\pi$ of $[a]$ set
\[
 C(\pi)=\sum_{t=1}^{a}w_{\pi_t}\prod_{s<t}r_{\pi_s}\prod_{s>t}f_{\pi_s}.
\]

\begin{lemma}[Adjacent exchange]
\label{lem:exchange}
Let $\pi$ and $\pi'$ agree except that $\pi$ places $i$ immediately before $j$ and $\pi'$
places $j$ immediately before $i$. Then
\[
 C(\pi)\le C(\pi')
 \iff
 w_id_j\le w_jd_i.
\]
\end{lemma}

\begin{proof}
Write $A$ for the product of $r$ over the slots preceding the pair and $B$ for the
product of $f$ over the slots following it. Both are common to $\pi$ and $\pi'$, as is
every summand indexed by a slot outside the pair. The two summands indexed by the pair
contribute $AB(w_if_j+r_iw_j)$ to $C(\pi)$ and $AB(w_jf_i+r_jw_i)$ to $C(\pi')$. Since
$A,B\ge0$, the comparison reduces to
\[
 w_if_j+r_iw_j\le w_jf_i+r_jw_i
 \iff
 w_i(f_j-r_j)\le w_j(f_i-r_i)
 \iff
 w_id_j\le w_jd_i.\qedhere
\]
\end{proof}

\begin{lemma}[The exchange relation is a total preorder]
\label{lem:preorder}
Assign to each slot the key
\[
 \kappa_i=\bigl(\mathrm{cls}(i),\,q_i\bigr),
 \qquad
 \mathrm{cls}(i)=
 \begin{cases}
 0,&d_i>0,\\
 1,&d_i=0,\\
 2,&d_i<0,
 \end{cases}
 \qquad
 q_i=
 \begin{cases}
 w_i/d_i,&d_i\ne0,\\
 0,&d_i=0,
 \end{cases}
\]
ordered lexicographically. Then $\kappa_i\le\kappa_j$ implies $w_id_j\le w_jd_i$, and the
lexicographic order is a total preorder on the slots.
\end{lemma}

\begin{proof}
Four cases, using $w_i,w_j\ge0$ throughout.

\emph{Same nonzero class.} If $d_i,d_j$ are both positive or both negative then
$d_id_j>0$, and dividing $w_id_j\le w_jd_i$ by $d_id_j$ shows it is equivalent to
$w_i/d_i\le w_j/d_j$, that is, to $q_i\le q_j$.

\emph{Positive before negative.} If $d_i>0$ and $d_j<0$ then $w_id_j\le0\le w_jd_i$, so
the inequality holds; here $\mathrm{cls}(i)=0<2=\mathrm{cls}(j)$.

\emph{Positive before zero.} If $d_i>0$ and $d_j=0$ then $w_id_j=0\le w_jd_i$, and
$\mathrm{cls}(i)=0<1=\mathrm{cls}(j)$.

\emph{Zero before negative.} If $d_i=0$ and $d_j<0$ then $w_id_j\le0=w_jd_i$, and
$\mathrm{cls}(i)=1<2=\mathrm{cls}(j)$.

Lexicographic order on pairs drawn from totally ordered sets is a total preorder.
\end{proof}

\begin{proof}[Proof of Theorem~\ref{thm:order}]
Let $\pi^\star$ sort the slots by nondecreasing $\kappa$; this is the order described in
the statement. Let $\pi$ be arbitrary. Because $\kappa$ induces a total preorder
(Lemma~\ref{lem:preorder}), $\pi$ can be transformed into $\pi^\star$ by a finite sequence
of adjacent transpositions, each swapping a pair that occurs in the order $j,i$ with
$\kappa_i\le\kappa_j$; this is the standard exchange-sorting argument, and it terminates
because each such transposition strictly decreases the number of pairs inverted with
respect to $\kappa$. By Lemmas~\ref{lem:preorder} and~\ref{lem:exchange}, no transposition
increases $C$. Hence $C(\pi^\star)\le C(\pi)$ for every $\pi$, so $\pi^\star$ is a global
minimiser. If $w_i=d_i=0$ then the inequality of Lemma~\ref{lem:exchange} holds in both
directions against every other slot, so such slots may be placed anywhere.
\end{proof}

\begin{remark}
\label{rem:orderfamily}
Theorem~\ref{thm:order} minimises $C$, which is one particular scalar bound. It does not
assert that the bound obtained with the optimal order is the smallest valid bound
obtainable by any means.
\end{remark}

\subsection{The state-resolved recursion and the pathwise expansion}

\begin{proof}[Proof of Theorem~\ref{thm:state}]
At a leaf all five stored quantities are exact: $\norm{F_v}=\norm{R_v}=\norm{z_v}$ and
the three error quantities vanish.

At an internal vertex, expand by Theorem~\ref{thm:subset} and split each child
discrepancy orthogonally as $D_i=D_i^{\parallel}+D_i^{\perp}$. Every resulting term is
$\mu_v(Z_1^\sigma,\ldots,Z_a^\sigma)$ for exactly one state vector
$\sigma\in\{R,\parallel,\perp\}^{a}$. Projecting the output onto the component $\chi$ and
applying the operator-norm inequality bounds the norm of that term by
$\beta_v^{\chi}(\sigma)\prod_i\norm{Z_i^{\sigma_i}}$, hence by
$\beta_v^{\chi}(\sigma)\prod_ib_{c_i}^{\sigma_i}$ once the inductive bounds are
substituted. The state vector $(R,\ldots,R)$ contributes the reduced value together with
the locally created residual: its projection onto $\ran P_v$ is $R_v$ itself, which is
subtracted, and its orthogonal component is $r_v(R_1,\ldots,R_a)$, of norm at most
$\rho_v\prod_ib_{c_i}^{R}$. This accounts for $\lambda_v^{\parallel}=0$ and
$\lambda_v^{\perp}\le\rho_v\prod_ib_{c_i}^{R}$. Summing over the remaining state vectors
and applying the triangle inequality gives the stated recursion. Taking, at any vertex,
the minimum of several independently valid upper bounds preserves validity.

For the count: there are $3^{a_v}-1$ state vectors other than $(R,\ldots,R)$ at $v$, and
each contributes one product of $a_v$ factors, so the number of combination steps at $v$
is $O(3^{a_v})$; summing over the vertices gives $O(\lvert T\rvert\,3^{a_{\max}})$.
\end{proof}

\begin{proof}[Proof of Corollary~\ref{cor:pathwise}]
Two separate assertions are proved.

\emph{Comparison.} We show $B_v\le\widehat B_v$ by induction on the subtree at $v$. At a
leaf both are $0$. At an internal vertex, the inductive hypothesis gives
$B_{c_j}\le\widehat B_{c_j}$ for every child, and the gains satisfy $h_{v,j}\ge0$, so
\[
 B_v
 \;\le\;\lambda_v+\sum_jh_{v,j}B_{c_j}
 \;\le\;\lambda_v+\sum_jh_{v,j}\widehat B_{c_j}
 \;=\;\widehat B_v,
\]
the first inequality being \eqref{eq:recurrence} and the last
Definition~\ref{def:majorant}. Nonnegativity of the gains is what makes the middle step
valid, and it holds because each $h_{v,j}$ is an operator norm.

\emph{Closed form.} We show
$\widehat B_v=\sum_{u\in\Int T_v}\lambda_u\prod_{e\in\mathrm{path}(u,v)}h_e$, where $T_v$
is the subtree rooted at $v$, again by induction. At a leaf both sides are $0$, the sum
being empty. At an internal vertex $v$, Definition~\ref{def:majorant} and the inductive
hypothesis give
\[
 \widehat B_v
 =\lambda_v+\sum_jh_{v,j}\!\!\sum_{u\in\Int T_{c_j}}\!\!\lambda_u
 \prod_{e\in\mathrm{path}(u,c_j)}h_e
 =\lambda_v+\sum_j\;\sum_{u\in\Int T_{c_j}}\!\!\lambda_u
 \prod_{e\in\mathrm{path}(u,v)}h_e ,
\]
because prefixing the edge from $c_j$ to $v$, whose gain is $h_{v,j}$, to the path from
$u$ to $c_j$ gives exactly the path from $u$ to $v$. The internal vertices of $T_v$ are
$v$ together with the disjoint union of the internal vertices of the subtrees $T_{c_j}$,
and the term for $u=v$ is $\lambda_v$ with an empty product. This is the asserted sum.
Taking $v=\varrho$ gives the statement, and combining the two assertions gives
$B_\varrho\le\widehat B_\varrho$.

\emph{Projected variant.} Setting $\lambda_\varrho=0$, which is
$\lambda^{\parallel}_\varrho=0$ from Theorem~\ref{thm:state}, removes the term $u=\varrho$
from the sum; the remaining derivation is unchanged, with the final edge of each path
carrying the projected-output gain.
\end{proof}

\begin{remark}
\label{rem:pathwisescope}
The closed form is an identity for $\widehat B$, and $\widehat B$ majorises $B$. It is not
an identity for $B$, and it is not an expansion of the multilinear cross terms of
Theorem~\ref{thm:subset}: the interaction terms indexed by subsets of size at least two
were absorbed into the gains $h_{v,j}$ when the scalar recurrence was formed. The result
attributes a share of the bound to individual sources, and not a share of the error.
\end{remark}

\subsection{Signed combinations and approximation of the maps}

\begin{proof}[Proof of Corollary~\ref{cor:signed}]
Linearity of $P$ and of the evaluation in each tree gives
$PF_{\mathcal F}-R_{\mathcal F}=\sum_\alpha c_\alpha(PF_{T_\alpha}-R_{T_\alpha})$, and the
triangle inequality bounds the norm by
$\sum_\alpha\lvert c_\alpha\rvert B^{\proj}_{T_\alpha}$. For the five-input ternary
associator, each of the two trees has two internal vertices, so
Theorem~\ref{thm:homogeneous} gives $B^{\proj}_{T_\alpha}\le(2-1)\rho ML=\rho ML$ and the
two terms give $2\rho ML$; the ambient bound is $2\cdot2\rho ML=4\rho ML$. Syntactically
identical trees may be combined exactly before this estimate is applied.
\end{proof}

\begin{lemma}[Sensitivity of the projected evaluation to the maps]
\label{lem:sensitivity}
Let $\{\mu_v\}$ and $\{\widehat\mu_v\}$ satisfy $\norm{\mu_v}_\op\le M$ and
$\norm{\mu_v-\widehat\mu_v}_\op\le\delta$ at every internal vertex, and let $R_\mu$ and
$R_{\widehat\mu}$ be the recursively projected evaluations of the two trees on the same
leaf data. Then for every subtree rooted at $v$,
\[
 \norm{R^{\mu}_v-R^{\widehat\mu}_v}\le k_v\,\delta\,(M+\delta)^{k_v-1}L_v
 \quad(k_v\ge1),
 \qquad
 \norm{R^{\mu}_v-R^{\widehat\mu}_v}=0
 \quad(k_v=0).
\]
\end{lemma}

\begin{proof}
At a leaf both evaluations equal $Q_{\tau(v)}z_v$, so the difference vanishes. At an
internal vertex $v$ of arity $a$, using $\norm{P_v}=1$ and adding and subtracting
$\widehat\mu_v(R^{\mu}_{c_1},\ldots,R^{\mu}_{c_a})$,
\begin{align*}
 \norm{R^{\mu}_v-R^{\widehat\mu}_v}
 &=\bigl\lVert P_v\bigl[\mu_v(R^{\mu}_{c_\bullet})-\widehat\mu_v(R^{\widehat\mu}_{c_\bullet})\bigr]\bigr\rVert\\
 &\le\underbrace{\bigl\lVert(\mu_v-\widehat\mu_v)(R^{\mu}_{c_1},\ldots,R^{\mu}_{c_a})\bigr\rVert}_{(\mathrm I)}
 +\underbrace{\bigl\lVert\widehat\mu_v(R^{\mu}_{c_\bullet})-\widehat\mu_v(R^{\widehat\mu}_{c_\bullet})\bigr\rVert}_{(\mathrm{II})}.
\end{align*}
For $(\mathrm I)$, Lemma~\ref{lem:magnitudes} gives
$\norm{R^{\mu}_{c_i}}\le M^{k_{c_i}}L_{c_i}$, so
$(\mathrm I)\le\delta M^{k_v-1}L_v\le\delta(M+\delta)^{k_v-1}L_v$.
For $(\mathrm{II})$, telescoping the arguments of $\widehat\mu_v$ one at a time and using
$\norm{\widehat\mu_v}_\op\le M+\delta$ together with the magnitude bounds
$\norm{R^{\mu}_{c_i}},\norm{R^{\widehat\mu}_{c_i}}\le(M+\delta)^{k_{c_i}}L_{c_i}$ from
Lemma~\ref{lem:magnitudes} applied with $M+\delta$ in place of $M$,
\[
 (\mathrm{II})
 \le(M+\delta)\sum_{j}\norm{R^{\mu}_{c_j}-R^{\widehat\mu}_{c_j}}
 \prod_{i\ne j}(M+\delta)^{k_{c_i}}L_{c_i}
 \le\Bigl(\sum_jk_{c_j}\Bigr)\delta(M+\delta)^{k_v-1}L_v,
\]
by the inductive hypothesis, the summand with $k_{c_j}=0$ vanishing. Since
$\sum_jk_{c_j}=k_v-1$, the two contributions add to $k_v\delta(M+\delta)^{k_v-1}L_v$.
\end{proof}

\begin{proof}[Proof of Proposition~\ref{prop:approximation}]
Insert the recursively projected \emph{exact} tree:
\[
 \norm{F_\mu-R_{\widehat\mu}}
 \le\norm{F_\mu-R_\mu}+\norm{R_\mu-R_{\widehat\mu}}.
\]
The first term is $E_T^{\amb}$ for the exact tree, bounded by $k\rho M^{k-1}L_T$ by
Theorem~\ref{thm:homogeneous}; the second is bounded by $k\delta(M+\delta)^{k-1}L_T$ by
Lemma~\ref{lem:sensitivity}. This gives the first estimate.

For the second, $R_\mu\in\ran P$ gives $PR_\mu=R_\mu$, so
\[
 \norm{PF_\mu-R_{\widehat\mu}}
 \le\norm{PF_\mu-R_\mu}+\norm{R_\mu-R_{\widehat\mu}},
\]
and the first term is $E_T^{\proj}$ for the exact tree, bounded by
$(k-1)\rho M^{k-1}L_T$ by Theorem~\ref{thm:homogeneous}.

Both estimates use only the declared hypotheses: $\norm{\mu_v}_\op\le M$ and
$\norm{r_v}_\op\le\rho$ for the exact maps, and $\norm{\mu_v-\widehat\mu_v}_\op\le\delta$.
In particular no bound is assumed on
$(I-P_v)\widehat\mu_v(P\,\cdot,\ldots,P\,\cdot)$; compare
Remark~\ref{rem:noapproxclosure}. No equality with the observed total error is asserted.
\end{proof}
